\ysrc{287}
然后如同 [Col10b] 中那样，最终通过进一步化归到下述显式描述而证得：即 Berger--Breuil [BB10] 所证明的 $\Gal_{\Q_p}$ 的二维晶体表示的 $p$-进局部 Langlands 对应。确立这一点之后，我们便有下述有用的推论。

\begin{cor}{7.3.2 (Dotto--E.--G.)}\label{thm:a4-7-3-2}
设 $p\ge3$。设 $A$ 是某个 $a\ge1$ 的有限型 $\OO/\varpi^a$-代数，且 $D_A$ 是秩 $2$、以 $A$ 为系数、行列式为 $\varepsilon$ 的投射 étale $(\varphi,\Gamma)$-模。则 $D_A^\natural\boxtimes\mathbf P^1$ 是有限 $A[[K]]$-模。
\end{cor}

\begin{proof}
这可由 $D_A^\natural\boxtimes\mathbf P^1$ 在 $\mathrm{GL}_2(\Q_p)$ 作用下的稳定性、$A[[K]]$ 的诺特性，以及下述事实推出：$D_A^\natural$ 作为格，是有限生成 $\mathbf A_A^+$-模。
\end{proof}

进而可以证明，$D_A^\natural\boxtimes\mathbf P^1$ 的构造与 $A$ 上的平坦基变换相容；我们利用这一点构造 $\XX$ 上带有连续 $\mathrm{GL}_2(\Q_p)$-作用的 pro-凝聚层 $L_\infty$。

\subsubsection*{7.3.3.\quad 对 Paškūnas 部分结果的重新解释}
现在回顾 [Paš13] 的一些主要结果，并用我们的函子 $\mathfrak A$ 加以重释。记 $(\mathrm{sm}\,G)^{\mathrm{l.adm}}$ 为 $\mathrm{sm}\,G$ 中由局部容许（locally admissible）表示组成的满子范畴。（按定义，表示 $\pi\in\mathrm{sm}\,G$ 是局部容许的，若每个向量 $v\in\pi$ 都光滑并生成一个容许表示；事实上对 $\mathrm{GL}_2(\Q_p)$，局部容许表示自动局部有限，从而对每个 $v$，由 $v$ 生成的表示都是有限长的。）

与 $\mathrm{sm}\,G$ 不同，范畴 $(\mathrm{sm}\,G)^{\mathrm{l.adm}}$ 容许块（block）分解，现解释如下。按定义，$(\mathrm{sm}\,G)^{\mathrm{l.adm}}$ 的一个块是（同构类的）不可约对象在如下等价关系下的一个等价类：该关系由 $\pi_1\sim\pi_2$（当 $\Ext^1_{\mathrm{sm}\,G}(\pi_1,\pi_2)\neq0$ 或 $\Ext^1_{\mathrm{sm}\,G}(\pi_2,\pi_1)\neq0$ 时）生成。

严格说来，本应写成「不可约容许对象」，但事实上 $\mathrm{sm}\,G$ 的不可约对象全都落在 $(\mathrm{sm}\,G)^{\mathrm{l.adm}}$ 中。确实，特征 $p$ 下 $\mathrm{GL}_2(\Q_p)$ 的、承认中心特征\footnote{后来 Berger 证明 [Ber12]：$\mathrm{GL}_2(\Q_p)$ 在 $\overline{\mathbb F}_p$ 上的不可约光滑表示必定承认中心特征。在我们的特定背景下，我们关心的是 $\mathrm{PGL}_2(\Q_p)$ 的表示，亦即中心特征平凡的 $\mathrm{GL}_2(\Q_p)$ 的表示，因此中心特征这一问题不会出现。}的绝对不可约光滑表示已由 Barthel--Livné 与 Breuil [BL94, Bre03a] 分类，特别地已被证明是容许的。非绝对不可约的情形很容易化归到绝对不可约情形（例如见 [Ber12, Prop. 1.2] 与 [Paš13, §5.13]）。

块 $\mathfrak B$ 如下：
\begin{enumerate}[label=(\arabic*),leftmargin=2.4em,itemsep=0.15em]
\item $\mathfrak B=\{\pi\}$，其中 $\pi$ 是中心特征平凡的超奇异不可约表示；
\item $\mathfrak B=\{\mathrm{Ind}_B^G(\chi\otimes\chi^{-1}),\mathrm{Ind}_B^G(\omega\chi^{-1}\otimes\omega^{-1}\chi)\}$，其中 $\chi:\Q_p^\times\to k^\times$ 是满足 $\chi^2\neq1,\omega^2$ 的某特征\tn{底本此处作 $\chi^2\neq1,\omega^2$；按上下文推断应读作 $\chi^2\notin\{1,\omega^2\}$（上下文推断，无外部依据），译文照录。}；
\item $\mathfrak B=\{\chi,\chi\otimes\mathrm{St}_G,\mathrm{Ind}_B^G(\omega\chi\otimes\omega^{-1}\chi)\}$，其中 $\chi:\Q_p^\times\to k^\times$ 是二次特征；以及
\item 不含绝对不可约对象的块。
\end{enumerate}

\ysrc{288}对我们而言，一个关键观察是：类型 (1)--(3) 的块与 $\XX(k)$ 之间存在双射。（为叙述方便，我们忽略类型 (4) 的块；这类块容易通过允许系数域 $k$ 的有限扩张来处理，它们确实对应于 $\XX$ 在 $k$ 的扩张上的点。）为描述该双射，回忆：对任一 Serre 权 $\sigma$，Hecke 代数
\[
\mathcal H(\sigma):=\End_G(c\text{-Ind}_K^G\sigma)
\]
等于 $k[T]$，其中 $T$ 是某个显式的 Hecke 算子；而绝对不可约表示都是表示
\[
(c\text{-Ind}_K^G\sigma)/(T-\lambda)
\]
（$\lambda\in k$）的子商。此外，除非 $\sigma=\sigma_{a,0}$ 或 $\sigma=\sigma_{a,p-1}$（对某个 $a$）且 $\lambda=\pm1$，这些表示都是（绝对）不可约的；在后一情形有两个 Jordan--Hölder 因子，其形式为 $\chi$、$\chi\otimes\mathrm{St}_G$（对某个 $\chi$）。

块与 $\XX(k)$ 之间的双射依赖于把 Hecke 算子 $T$ 看作 $\XX$ 的分支 $\XX(\sigma|\sigma^{co})$ 中 $\mathbf A^1$ 上的坐标。具体做法如下，仍如上面那样限于一般情形叙述。

\begin{defn}{7.3.4}\label{thm:a4-7-3-4}
设 $\sigma$ 是一般 Serre 权。上文已定义态射 $\XX(\sigma|\sigma^{co})\to\XX(\sigma|\sigma^{co})\subset\XX$，它由 $\pi_\sigma$ 或 $\pi_{\sigma^{co}}$ 之一给出；并固定了 $\XX(\sigma|\sigma^{co})$（同构于 $\mathbf P^1$）上的坐标 $T$。若该态射由 $\pi_\sigma$ 给出，则定义 $f_\sigma:\Spec\mathcal H(\sigma)\to\XX$ 为 $T\mapsto T$；若该态射由 $\pi_{\sigma^{co}}$ 给出，则令 $T\mapsto T^{-1}$。
\end{defn}

于是若 $x$ 是 $\XX$ 的 $k$-点，相应的块由
\[
\mathfrak B_x=\bigcup_{\sigma\ \text{s.t.}\ x\in f_\sigma(\mathbf A^1)}\mathrm{JH}\Bigl(c\text{-Ind}_K^G(\sigma)\otimes_{\mathcal H(\sigma),x}k\Bigr)
\]
给出。更确切地说，$\mathfrak B_x$ 是 $(\mathrm{sm}\,G)_k$ 的一个块，而映射 $x\mapsto\mathfrak B_x$（适当延拓到非一般 Serre 权的情形）是从 $\XX(k)$ 到「含绝对不可约表示的 $\mathrm{sm}\,G$ 的块」之集合的双射。

若 $\mathfrak B$ 是 $(\mathrm{sm}\,G)^{\mathrm{l.adm}}$ 的块，则记 $(\mathrm{sm}\,G)^{\mathrm{l.adm}}_{\mathfrak B}$ 为 $(\mathrm{sm}\,G)^{\mathrm{l.adm}}$ 中由「所有不可约子商都落在给定块 $\mathfrak B$ 中」的表示组成的满子范畴。于是有块分解
\[
(\mathrm{sm}\,G)^{\mathrm{l.adm}}\isom\prod_{\mathfrak B}(\mathrm{sm}\,G)^{\mathrm{l.adm}}_{\mathfrak B}.
\]

在 [Paš13] 中，Paškūnas 建立了范畴 $(\mathrm{sm}\,G)^{\mathrm{l.adm}}_{\mathfrak B}$ 与若干伪形变环上的模范畴之间的范畴等价。我们把 Paškūnas 的结果（及其证明）用作定理 7.2.1 证明的输入；但为此目的，必须把这些结果（某些情形下还要稍加推广）重新表述为关于「映入叠上层范畴的全忠实函子」的陈述。

设 $\bar\theta:\Gal_{\Q_p}\to\mathrm{GL}_2(k)$ 是半单表示（等价地，伪表示），对应于一点 $x\in\XX(k)$。假设（为叙述简便）相应的块 $\mathfrak B_x$ 属于上述情形 (1) 或 (2)；也就是说，假设 $x$ 不是 $\XX(\sigma|\sigma^{co})$ 中满足 $T=\pm1$ 的点，其中 $\{\sigma,\sigma^{co}\}$ 形如 $\{\sigma_{a,0},\sigma_{a+1,p-3}\}$。记
\[
(\mathrm{sm}\,G)^{\mathrm{l.adm}}_{\bar\theta}:=(\mathrm{sm}\,G)^{\mathrm{l.adm}}_{\mathfrak B_x}.
\]

\ysrc{289}
以 $\XX_{\bar\theta}$ 记 Carl Wang-Erickson 的与 $\bar\theta$ 相伴的伽罗瓦表示叠。按定义，$\XX_{\bar\theta}$ 由下述性质刻画：若 $A$ 是 $p$ 在其中幂零的 $\OO$-代数，则 $\XX_{\bar\theta}(A)$ 是连续态射
\[
\rho:\Gal_{\Q_p}\to\mathrm{GL}_2(A)
\]
的群胚，其中 $\det\rho=\varepsilon^{-1}$，且 $\rho$ 模 $\varpi$ 相伴的伪表示等于 $\bar\theta$（此处 $A$ 取离散拓扑，$\Gal_{\Q_p}$ 取自然的投射有限拓扑）。

Wang-Erickson 在 [WE18] 中证明 $\XX_{\bar\theta}$ 是形式代数叠；又由 [EG23, Thm. 6.7.2] 与 [EG24, Thm. 10.2.2]，存在自然的单态射 $\XX_{\bar\theta}\into\XX$，它在底拓扑空间上诱导闭浸入（其像对应于那些半单化同构于 $\bar\theta$ 的表示 $\bar\rho:\Gal_{\Q_p}\to\mathrm{GL}_2(\mathbb F_p)$）。该态射在 $\XX_{\bar\theta}$ 的有限型点处是 versal 的，因此我们自然预期它是形式代数叠 $\XX$ 沿 $\XX_{\mathrm{red}}$ 的闭子叠 $\ZZ_{\bar\theta}$ 的某种完备化；$\ZZ_{\bar\theta}$ 的 $\mathbb F_p$-点正是那些半单化为 $\bar\theta$ 的 $\bar\rho$。

这一预期可用凝聚完备性（coherent completeness）[AHR20, Defn. 2.1] 的语言精确表述如下。以 $\mathcal I_{\bar\theta}$ 记 $\ZZ_{\bar\theta}$ 在 $\XX$ 中的理想层，以 $\XX^{[n]}_{\bar\theta}$ 记 $\XX$ 中由 $\mathcal I^{n+1}_{\bar\theta}$ 截出的闭子叠。则可以证明：对每个 $n$，$\XX^{[n]}_{\bar\theta}$ 是 $\XX_{\bar\theta}$ 的闭子叠；且自然函子（由限制给出）
\[
\Coh(\XX_{\bar\theta})\to\varprojlim_n\Coh(\XX^{[n]}_{\bar\theta})
\]
是范畴等价；也就是说，$\XX_{\bar\theta}$ 上的凝聚层等同于 $\XX$ 中 $\ZZ_{\bar\theta}$ 的各无穷小邻域上的一族相容凝聚层。

给定上述构造，下述定理本质上只是拆解 [Paš13] 主要结果的部分证明。以 $(\mathrm{sm}\,G)^{\mathrm{fg.adm}}_{\bar\theta}$ 记 $(\mathrm{sm}\,G)^{\mathrm{l.adm}}_{\bar\theta}$ 中由有限生成表示组成的满子范畴。

\begin{thm}{7.3.5 (Paškūnas$+\varepsilon$)}\label{thm:a4-7-3-5}
设如上有伪表示 $\bar\theta:\Gal_{\Q_p}\to\mathrm{GL}_2(k)$，其对应的块落在情形 (1) 或 (2)。则函子
\begin{equation}\tag{7.3.6}\label{eq:a4-7-3-6}
(L_\infty)|_{\XX_{\bar\theta}}\otimes_{\OO[[G]]}-
\end{equation}
是从 $(\mathrm{sm}\,G)^{\mathrm{fg.adm}}_{\bar\theta}$ 到 $\XX_{\bar\theta}$ 上凝聚层的范畴的正合全忠实函子。
\end{thm}

\begin{proof}
以 $R^{\mathrm{ps}}_{\bar\theta}$ 记完备诺特局部 $\OO$-代数，它参数化 $\bar\theta$ 在有限型 Artin 局部 $\OO$-代数上的、行列式固定为 $\varepsilon^{-1}$ 的形变；以 $\theta^u$ 记 $\bar\theta$ 在 $R^{\mathrm{ps}}_{\bar\theta}$ 上的泛形变。我们以如下方式定义 Cayley--Hamilton 代数 $\widetilde R_{\bar\theta}$：以 $J$ 记 $R^{\mathrm{ps}}_{\bar\theta}[[\Gal_{\Q_p}]]$ 中由元素 $g^2-\theta^u(g)g+(\varepsilon^{-1})(g)$（$g\in\Gal_{\Q_p}$）生成的双边理想的闭包，然后令
\[
\widetilde R_{\bar\theta}:=R^{\mathrm{ps}}_{\bar\theta}[[\Gal_{\Q_p}]]/J.
\]
\ysrc{290}
如同第 3 节，以 $V$ 记 Colmez 的 Montreal 函子。Paškūnas 在 [Paš13] 中证明环 $\widetilde R_{\bar\theta}$ 作为 $R^{\mathrm{ps}}_{\bar\theta}$-模是秩 $4$ 的自由模；而 [Paš13] 的结果可以解释为：函子 $\pi\mapsto V(\pi)$ 在 $(\mathrm{sm}\,G)^{\mathrm{fg.adm}}_{\bar\theta}$ 与有限长 $\widetilde R_{\bar\theta}$-模范畴之间诱导等价。

存在泛二维 $\Gal_{\Q_p}$-表示 $\mathcal V$ 铺在 $\XX_{\bar\theta}$ 上，可看作带 $\widetilde R_{\bar\theta}$ 作用的秩 $2$ 局部自由 $\OO_{\XX_{\bar\theta}}$-模层。于是利用 [Col10b] 中建立的 $D^\natural\boxtimes\mathbf P^1$ 与 Montreal 函子 $V$ 之间的关系，可以证明函子 (7.3.6) 自然等价于函子
\[
\mathcal V\otimes_{\widetilde R_{\bar\theta}}V(-),
\]
于是由 [Paš13] 中证明的定理容易推出结论。
\end{proof}

类型 (3) 的块（即 $x$ 是 $\XX(\sigma|\sigma^{co})$ 中满足 $T=\pm1$ 的点，其中 $\{\sigma,\sigma^{co}\}$ 形如 $\{\sigma_{a,0},\sigma_{a+1,p-3}\}$ 的情形）更为微妙，因为与其他块不同，我们会看到导出（derived）现象。这一情形最近由 Christian Johansson、James Newton 与 Carl Wang-Erickson [JNWE24] 研究，他们的结果可用于证明下述定理。

\begin{thm}{7.3.7 (Johansson--Newton--Wang-Erickson$+\varepsilon$)}\label{thm:a4-7-3-7}
设如上有伪表示 $\bar\theta:\Gal_{\Q_p}\to\mathrm{GL}_2(k)$，其对应的块 $\mathfrak B_{\bar\theta}=\mathfrak B_x$ 落在情形 (3)。则函子
\begin{equation}\tag{7.3.8}\label{eq:a4-7-3-8}
(L_\infty)|_{\XX_{\bar\theta}}\otimes^{\mathbf L}_{\OO[[G]]}-
\end{equation}
是从 $D(\mathfrak B^{\mathrm{fg}}_{\bar\theta})$ 到 $\XX_{\bar\theta}$ 上凝聚模的有界导出范畴的、幅度为 $[-1,0]$ 的全忠实函子。
\end{thm}

\subsubsection*{7.3.9.\quad 一个有用的消解}
设 $N$ 是通常（上三角）Iwahori 子群 $I$ 的正规化子，$\delta$ 是 $N/ZI$ 的非平凡二次特征。若 $\pi$ 是任一光滑 $G$-表示，则有 $\OO[[G]]$-模复形
\begin{equation}\tag{7.3.10}\label{eq:a4-7-3-10}
0\to c\text{-Ind}_N^G(\delta\pi)\to c\text{-Ind}_K^G(\pi)\to\pi\to0,
\end{equation}
其中第三个箭头是显然的满射 $c\text{-Ind}_K^G\pi\to\pi$。第二个箭头的定义如下。记 $\Pi=\begin{pmatrix}0&1\\p&0\end{pmatrix}$。则态射
\[
v\mapsto\Pi\otimes\Pi^{-1}v-1\otimes v\in\OO[G]\otimes_{\OO[K]}\pi=c\text{-Ind}_K^G\pi
\]
是 $N$-等变态射 $\delta\pi\to c\text{-Ind}_K^G\pi$，它诱导第二个箭头 $c\text{-Ind}_N^G\delta\pi\to c\text{-Ind}_K^G\pi$。容易证明复形 (7.3.10) 无环。

\begin{rem}{7.3.11}\label{thm:a4-7-3-11}
消解 (7.3.10) 对任意 $F/\Q_p$ 的 $\mathrm{PGL}_2(F)$ 都成立，可用于把关于一般 $\pi$ 的某些断言化归到 $\pi=c\text{-Ind}_K^G\sigma$（$\sigma$ 为 Serre 权）的情形。特别地，可用如下方式把连续函子（例如我们假设的函子 $\mathfrak A$）在一般 $\pi$ 上的求值化归到在 $c\text{-Ind}_K^G\sigma$ 上的求值。由于我们假设 $p>2$，$c\text{-Ind}_N^G\delta\pi$ 是 $c\text{-Ind}_I^G\pi=c\text{-Ind}_K^G(c\text{-Ind}_I^K\pi)$ 的直和项，故可化归到研究形如 $c\text{-Ind}_K^GV$ 的表示。
\ysrc{291}由于紧诱导与取滤过余极限相容，我们随后化归到 $V$ 有限生成的情形。由于紧诱导正合，我们进一步化归到 $\pi$ 形如 $c\text{-Ind}_K^G\sigma$（$\sigma$ 为某 Serre 权）的情形，如所断言。
\end{rem}

\subsubsection*{7.3.12.\quad 函子的定义}
现在可以定义我们的函子了。既然已在 $\XX$ 上定义了 $\OO[[G]]$-表示的 pro-凝聚层 $L_\infty$，我们自然想通过
\begin{equation}\tag{7.3.13}\label{eq:a4-7-3-13}
\pi\mapsto L_\infty\otimes^{\mathbf L}_{\OO[[G]]}\pi
\end{equation}
来定义它。然而这一先验地并不显然性质良好，故我们分阶段进行。首先，若 $\pi$ 形如 $c\text{-Ind}_K^GV$ 且 $V$ 有限生成，则注意
\[
L_\infty\otimes^{\mathbf L}_{\OO[[G]]}c\text{-Ind}_K^GV=L_\infty\otimes^{\mathbf L}_{\OO[[K]]}V,
\]
利用推论 7.3.2 可证它是凝聚 $\OO_{\XX}$-模的上有界复形。我们断言它其实是集中于 $0$ 次的层。为看出这点，只需（对上同调次数作降序归纳，并利用上同调群的凝聚性）在每个子叠 $\XX_{\bar\theta}$ 上限制后验证，而那里结论由定理 7.3.5 给出。

过渡到余极限，我们看到对任意 $V$（无论有限生成与否），$L_\infty\otimes^{\mathbf L}_{\OO[[G]]}c\text{-Ind}_K^GV$ 是 $0$ 次拟凝聚层。由注 7.3.11，推出对任意光滑表示 $\pi$，$L_\infty\otimes^{\mathbf L}_{\OO[[G]]}\pi$ 是集中于次数 $[-1,0]$ 的拟凝聚层复形；且若 $\pi$ 有限表现（等价地，有限生成），则该复形具有凝聚上同调层。我们断定 $L_\infty\otimes^{\mathbf L}_{\OO[[G]]}-$ 定义了幅度落在 $[-1,0]$ 中的函子，并把有限表现表示的有界复形送到 $\OO_{\XX}$-模的有界凝聚复形。

\subsubsection*{7.3.14.\quad 局部化与粘合}
若 $\sigma$ 是 Serre 权，我们通过定义 7.3.4 启发式地把 $c\text{-Ind}_K^G\sigma$ 看作「生活」在 $\XX$ 中 $\mathbf A^1=\Spec k[T]=\Spec\mathcal H(\sigma)$ 的某个拷贝之上。这一直观是 [DEG23] 中发展的 $\mathrm{sm}\,G$ 局部化理论的基础；现解释如何使之精确。本节所述结果及其证明大体纯属表示论（偶尔用到块分类，后者在 [Paš13] 中用 $p$-进局部 Langlands 对应证明），当然其动机来自几何图景。

若 $Y$ 是 $X$ 的闭子集\tn{底本对模叠记 $\XX$，对其底拓扑空间则时而记作 $X$（如本句及「$|X|$ 的开子集」）；二者含义不同，译文照录。}，以 $(\mathrm{sm}\,G)_Y$ 记 $\mathrm{sm}\,G$ 中由「所有不可约子商都落在对应于 $Y$ 的闭点的块中」的表示组成的子范畴。若 $U$ 是 $X$ 的开子集，则记
\[
(\mathrm{sm}\,G)_U:=\mathrm{sm}\,G/(\mathrm{sm}\,G)_Y,
\]
其中 $Y:=X\setminus U$。不难验证子范畴 $(\mathrm{sm}\,G)_Y$ 按通常意义是局部化子范畴：即典范商函子 $(j_U)^*:\mathrm{sm}\,G\to(\mathrm{sm}\,G)_U$ 有全忠实右伴随
\[
(j_U)_*:(\mathrm{sm}\,G)_U\to\mathrm{sm}\,G.
\]
在 $U$ 明确时，记 $j_*$ 为 $(j_U)_*$，$j^*$ 为 $(j_U)^*$。

下述结果即 [DEG23, Prop. 3.1.7]，它证实了关于函子 $j_*$、$j^*$ 应有行为的某些基本直观。

\ysrc{292}
\begin{prop}{7.3.15}\label{thm:a4-7-3-15}
设 $Y$ 是 $X$ 的闭子集，记 $U:=X\setminus Y$。
\begin{enumerate}[label=(\arabic*),leftmargin=2.4em,itemsep=0.15em]
\item 若 $\sigma$ 是 Serre 权，且对某个 $g\in\mathcal H(\sigma)$ 有 $f_\sigma^{-1}(Y)=V(g)$（$\Spec\mathcal H(\sigma)$ 的闭子集），则自然映射
\[
c\text{-Ind}_K^G\sigma\to(c\text{-Ind}_K^G\sigma)[1/g]
\]
可等同于态射
\[
c\text{-Ind}_K^G\sigma\to j_*j^*c\text{-Ind}_K^G\sigma.
\]
\item 函子 $j_*$ 正合，且与滤过余极限交换。
\end{enumerate}
\end{prop}

若 $\{U_0,\dots,U_n\}$ 是 $X$ 的任一有限开覆盖，则对 $\mathrm{sm}\,G$ 的任一对象 $\pi$，得到函子性的 Čech 消解
\begin{equation}\tag{7.3.16}\label{eq:a4-7-3-16}
0\to\pi\to\prod_i(j_i)_*(j_i)^*\pi\to\cdots\to(j_{0,\dots,n})_*(j_{0,\dots,n})^*\pi\to0,
\end{equation}
其中照例以 $U_{i\dots k}$ 记 $U_i\cap\cdots\cap U_k$，以 $j_{i\dots k}$ 记 $j_{U_{i\dots k}}$，微分由通常公式给出。下述结果即 [DEG23, Prop. 3.2.3]。

\begin{prop}{7.3.17}\label{thm:a4-7-3-17}
对 $\mathrm{sm}\,G$ 的任一对象 $\pi$，消解 (7.3.16) 无环。
\end{prop}

\begin{proof}
利用消解 (7.3.10)，如注 7.3.11 那样化归到 $\pi=c\text{-Ind}_K^G\sigma$（$\sigma$ 为 Serre 权）的情形。此时命题 7.3.15 允许我们化归到标准断言：仿射概形按区分开集的有限开覆盖所对应的 Čech 复形无环。
\end{proof}

由命题 7.3.17 推出，$\mathrm{sm}\,G$ 是 $X$ 上阿贝尔范畴的叠，如定理 7.2.3 所断言。

我们需要一个关于表示范畴的粘合陈述，它对应于层的 Beauville--Laszlo 粘合。记 $(\mathrm{sm}\,G)^{\mathrm{fg}}$ 为 $\mathrm{sm}\,G$ 中由有限生成对象组成的满子范畴。类似地，对 $X$ 的任一闭子集 $Y$，记 $(\mathrm{sm}\,G)^{\mathrm{fg}}_Y$ 为 $(\mathrm{sm}\,G)_Y$ 中由有限生成对象组成的满子范畴。可以构造形式 pro-对象范畴 $\mathrm{Pro}((\mathrm{sm}\,G)^{\mathrm{fg}}_Y)$。

\begin{defn}{7.3.18}\label{thm:a4-7-3-18}
若 $\pi$ 是 $(\mathrm{sm}\,G)^{\mathrm{fg}}$ 的对象，$Y$ 是 $X$ 的闭子集，则记 $\hat\pi_Y$ 为 $\mathrm{Pro}((\mathrm{sm}\,G)^{\mathrm{fg}}_Y)$ 中的对象「$\varprojlim_I$」$\pi'$，其中 $\pi'$ 遍历 $\pi$ 的、落在 $(\mathrm{sm}\,G)^{\mathrm{fg}}_Y$ 中的所有商。
\end{defn}

下述结果即 [DEG23, Thm. 3.8.1]。

\begin{thm}{7.3.19}\label{thm:a4-7-3-19}
若 $Y$ 有限，则函子 $\pi\mapsto\hat\pi_Y$ 正合，且由沿 $Y$ 的完备化与 $U$ 上的局部化诱导的典范函子
\[
\begin{aligned}
(\mathrm{sm}\,G)^{\mathrm{fg}}\longrightarrow{}&\mathrm{Pro}\bigl((\mathrm{sm}\,G)^{\mathrm{fg}}_Y\bigr)\\
&\times_{\mathrm{Pro}((\mathrm{sm}\,G)^{\mathrm{fg}}_Y)_U}(\mathrm{sm}\,G)^{\mathrm{fg}}_U
\end{aligned}
\]
是范畴等价。
\end{thm}

\ysrc{293}
\subsubsection*{7.3.20.\quad $\mathfrak A$ 的一般描述}
给定 $|X|$ 的开子集 $U$，有局部化范畴 $(\mathrm{sm}\,G)_U$，以及 $\XX$ 的开子叠 $\mathcal U:=\pi^{-1}(U)$。现在取 $U$ 为有限集 $Y$ 的补集，$Y$ 由对应于类型 (1) 或 (3) 的块的点组成；特别地我们删去了标记点。则 $U$ 在 $X$ 中稠密；如下将解释，实际上可以显式描述 $\mathfrak A$ 在 $(\mathrm{sm}\,G)_U$ 上的限制，并证明它是映入 $D(\mathcal U)$ 的全忠实函子。粗略的想法是循 [Paš13] 对类型 (2) 的块的策略，但「沿非分歧方向代数化」。

开集 $U$ 是它与 $|X|$ 各不可约分支之交的不交并，我们逐个处理这些交。本讨论中只限于描述 $U\cap X(\sigma|\sigma^{co})$ 的情形，其中 $\sigma,\sigma^{co}$ 是一对相伴的、非常一般的 Serre 权。该交在 §7.2.22 中记作 $U(\sigma|\sigma^{co})$；为简化记号，现在重新定义 $U$ 为 $U:=U(\sigma|\sigma^{co})$。$U$ 在 $\XX$ 中的原像为开子叠 $\mathcal U:=\mathcal U(\sigma|\sigma^{co})$。

于是我们的目标是描述 $\mathfrak A$ 与到 $\mathcal U$ 的限制之复合；结果表明它经过 $(\mathrm{sm}\,G)_U$ 分解。我们对该复合的描述将是 Morita 理论的，用 $(\mathrm{sm}\,G)_U$ 的一个投射 pro-生成元来表述；而用我们定义 $\mathfrak A$ 所用的层 $L_\infty$ 来构造这个 pro-生成元，是最简单的办法。事实上，使用 $L_\infty$ 的对偶版本 $L^\vee_\infty$ 更方便，它的构造方式与 $L_\infty$ 本身类似（另见下文 §7.3.22）。

现在写下
\[
\mathcal U\setminus\XX(\sigma^{co})=[(\Spf\OO[s,t^{\pm1},x,y^{\pm1}]^\wedge/\widehat{\mathbb G}_m]=[\Spf A/\mu_2],
\]
其中 $A=\OO[s,t^{\pm1},x]^\wedge$（帽子表示 $(p,s,x)$-进完备化），$\mu_2$ 的作用平凡。\tn{底本此式作「$[(\Spf\OO[s,t^{\pm1},x,y^{\pm1}]^\wedge/\widehat{\mathbb G}_m]$」，左圆括号无配对；译文照录。}可以把 $L^\vee_\infty|_{\mathcal U\setminus\XX(\sigma^{co})}$ 看作带 $\mu_2$ 作用的拓扑 $A[[G]]$-模 $P_\alpha$。经过一些工作并与 [Paš13] 的结果比较，可以证明 $P_\alpha$ 是 $\mathrm{Pro}((\mathrm{sm}\,G)_U)$ 的「pro-投射」对象，意即：若 $\pi\to\pi'$ 是 $(\mathrm{sm}\,G)_U$ 中的满态射，则 $\Hom(P_\alpha,\pi)\to\Hom(P_\alpha,\pi')$ 满射。（由于阿贝尔群满射的投射极限未必满射，这并不能推出 $P_\alpha$ 在 $\mathrm{Pro}((\mathrm{sm}\,G)_U)$ 中投射。）

交换 $\sigma$ 与 $\sigma^{co}$ 的角色，得到另一个 pro-投射对象 $P_\beta$。$P_\alpha$ 与 $P_\beta$ 共同生成 $(\mathrm{sm}\,G)_U$。利用定理 7.3.5 计算它们在闭点处的完备化，可以证明 $\mathfrak A(P_\alpha)$ 与 $\mathfrak A(P_\beta)$ 是 $\mathcal U$ 上某些显式的线丛。事实上，在适当解释记号后，它们就是例 C.4.7 中所考虑的分次模 $B(1)$ 与 $B(-1)$ 对应的线丛；而函子 $\mathfrak A:(\mathrm{sm}\,G)_U\to D(\mathcal U)$ 的全忠实性是命题 C.4.9 的直接推论。

\subsubsection*{7.3.21.\quad 函子的全忠实性}
如 §7.3.20 开头那样，以 $U$ 记有限集 $Y$ 的补集，$Y$ 由对应于类型 (1) 或 (3) 的块的点组成。如同往常忽略类型 (3) 的块的问题（它们需要一些额外论证），$\mathfrak A$ 的全忠实性可由 $\mathfrak A:(\mathrm{sm}\,G)_U\to D(\mathcal U)$ 的全忠实性、连同定理 7.3.19 与定理 7.3.5 推出。

\subsubsection*{7.3.22.\quad 与对偶的相容性}
对任一有限型 $\OO/\varpi^a$-代数 $A$，存在 $\mathrm{GL}_2(\Q_p)$-表示的显然短正合列
\begin{equation}\tag{7.3.23}\label{eq:a4-7-3-23}
0\to D_A^\natural\boxtimes\mathbf P^1\to D_A\boxtimes\mathbf P^1\to(D_A\boxtimes\mathbf P^1)/(D_A^\natural\boxtimes\mathbf P^1)\to0.
\end{equation}

\ysrc{294}
该短正合列在 [Col10b] 中起重要作用；利用 Colmez 的结果（通过化归到 $A$ 为 Artin 代数的情形），可以证明 $D_A^\natural\boxtimes\mathbf P^1$ 与 $(D_A\boxtimes\mathbf P^1)/(D_A^\natural\boxtimes\mathbf P^1)$ 之间存在完美对偶配对（最终来自 Tate 局部对偶的留数配对）。

以 $\mathbf D:D^b_{\mathrm{f.p.}}(\mathrm{sm}\,G)\to D^b_{\mathrm{f.p.}}(\mathrm{sm}\,G)$ 记由
\[
\mathbf D(-)=R\Hom(-,\OO[[G]])[4]
\]
定义的反变函子。利用消解 (7.3.10) 容易验证 $\mathbf D$ 的幅度为 $[0,1]$。我们还以 $\mathbf D_{\XX}$ 记由 Grothendieck--Serre 对偶给出的 $D_{\mathrm{coh}}(\XX)$ 的反等价，即 $R\Hom_{\OO_{\XX}}(-,\omega_{\XX})$，其中 $\omega_{\XX}$ 是 $\XX$ 的对偶层，置于 $0$ 次。

由 $L_\infty$ 的构造，可用 (7.3.23) 构造一个扩张
\[
0\to L_\infty\to?\to\mathbf D_{\XX}L_\infty\to0.
\]
该扩张类属于
\[
\Hom_{\OO[[G]]}\bigl(\mathbf D_{\XX}L_\infty,L_\infty[1]\bigr)=\Ext^1_{\OO[[G]]}\bigl(\mathbf D_{\XX}L_\infty,L_\infty\bigr)[1],
\]
而与该扩张类复合，给出 $D^b_{\mathrm{f.p.}}(\mathrm{sm}\,G)$ 上函子的自然态射
\[
R\Hom_{\OO[[G]]}(-,\mathbf D_{\XX}L_\infty)\to R\Hom_{\OO[[G]]}(-,L_\infty)[1];
\]
我们预期可以证明它是同构，从而得到所预期的与对偶的相容性
\[
\mathfrak A\circ\mathbf D\isom(\mathbf D_{\XX}\circ\mathfrak A)[3].
\]
如注 6.1.35 所述，这意味着各种层
\[
\mathfrak A\Bigl(\bigl(c\text{-Ind}_K^GW_\lambda\otimes\sigma^{\mathrm{crys},\circ}(\tau)\bigr)^\wedge\Bigr)
\]
都是 Cohen--Macaulay 的。由于这一性质可通过在 versal 环上拉回来验证，它也更直接地由函子 $\mathfrak A$ 的构造、它与 §7.3.3 所讨论的 Paškūnas 结果的关系、以及 [Paš15] 的结果推出。

\subsection{$\mathrm{GL}_2(\Q_p)$ 的 Banach 情形——III.\ 例}
现在描述函子 $\mathfrak A$ 的一些显式取值，从其在 Serre 权的紧诱导与（绝对）不可约表示上的值开始，然后转向某些与驯顺潜在 Barsotti--Tate 伽罗瓦表示有关的情形。后一类例子还说明了 §6.2.11 中所假定的 Banach 情形与解析情形之间的相容性。

\subsubsection*{7.4.1.\quad Serre 权的紧诱导}
考虑 Serre 权 $\sigma_{a,b}$。首先假设 $b\neq p-1$，即 $\sigma_{a,b}$ 一般。则 $\mathfrak A(c\text{-Ind}_K^G\sigma_{a,b})$ 是 $\XX(\sigma_{a,b})$ 上的线丛。这可预期，例如通过比较相应的 Taylor--Wiles 拼接模即可看出：由 Fontaine--Laffaille 理论，它们在支撑上是自由模；但究竟得到 $\XX(\sigma_{a,b})$ 上哪个线丛，则不那么显然。为描述它，注意由命题 7.2.11（以及 $\mathbf A^2\setminus\{0\}$ 的 Picard 群平凡），有 $\mathrm{Pic}(\XX(\sigma_{a,b}))=\mathbf Z$，且可验证 $\mathfrak A(c\text{-Ind}_K^G\sigma_{a,b})$ 是对应于 $-1\in\mathbf Z$ 的线丛。（这其实由与对偶的相容性所迫；注 6.1.33 的讨论表明 $\mathfrak A(c\text{-Ind}_K^G\sigma_{a,b})$ 必是 $\XX(\sigma_{a,b})$ 上唯一的自对偶线丛。关于 $\mathrm{GL}_2(\Q_{p^2})$ 的相关计算另见 §7.7。）

若 $b=p-1$，则 $\mathfrak A(c\text{-Ind}_K^G\sigma_{a,p-1})$ 是并
\begin{equation}\tag{7.4.2}\label{eq:a4-7-4-2}
\XX(\sigma_{a,0})\cup\XX(\sigma_{a,p-1})^+\cup\XX(\sigma_{a,p-1})^-
\end{equation}
上的线丛。

\ysrc{295}
我们预期该线丛可如下显式描述：由命题 7.2.13，叠 $\XX(\sigma_{a,p-1})^\pm$ 的 Picard 群也同构于 $\mathbf Z$；我们预期可验证 $\mathfrak A(c\text{-Ind}_K^G\sigma_{a,p-1})$ 在 (7.4.2) 中各不可约分支上的限制都对应于 $-1\in\mathbf Z$。

\subsubsection*{7.4.3.\quad 绝对不可约表示}
§7.4.1 的结果使我们能计算 $\mathfrak A$ 在绝对不可约表示上的值。如 §7.3.3 所讨论，它们是表示
\[
(c\text{-Ind}_K^G\sigma)/(T-\lambda)
\]
（$\lambda\in k$）的子商；且除非 $\sigma=\sigma_{a,0}$ 或 $\sigma=\sigma_{a,p-1}$（$a$ 为某值）且 $\lambda=\pm1$，这些表示事实上绝对不可约。假设不在此情形，则可由 $T-\lambda$ 在线丛 $\mathfrak A(c\text{-Ind}_K^G\sigma)$ 上的单射性推出：$\mathfrak A((c\text{-Ind}_K^G\sigma)/(T-\lambda))$ 是支撑在闭轨迹 $\pi_\sigma^{-1}(\lambda)\cap\XX(\sigma)$ 上的秩 $1$ 层。（当 $\lambda\neq0$ 时，我们将在下文 §7.4.4 中更精确地描述这个秩 $1$ 层。）

不出所料，平凡表示与 Steinberg 表示表现出更复杂的行为，这映照了特征零情形下 $\ell$-进表示的相应图景（见注 C.4.5）。存在 $G$-表示的短正合列
\[
0\to c\text{-Ind}_K^G\sigma_{0,0}\to c\text{-Ind}_K^G\sigma_{0,p-1}\to\mathrm{St}\oplus(\mathrm{nr}_{-1}\circ\det)\otimes\mathrm{St}\to0
\]
与
\[
0\to c\text{-Ind}_K^G\sigma_{0,p-1}\to c\text{-Ind}_K^G\sigma_{0,0}\to\mathbf 1\oplus\mathrm{nr}_{-1}\circ\det\to0,
\]
它们与上述 $\mathfrak A(c\text{-Ind}_K^G\sigma_{a,b})$ 的描述相结合，可用于证明 $\mathfrak A(\mathrm{St})$ 是 $\XX(\mathrm{Sym}^{p-1})^+$ 上的线丛，而 $\mathfrak A(\mathbf 1)$ 是 $\XX(\mathrm{Sym}^{p-1})^+$ 上另一个线丛经过上同调次数 $-1$ 的平移。

\subsubsection*{7.4.4.\quad 类型 (2) 的块}
更详细地考虑类型 (2) 的块的情形。此时 $\mathfrak B=\{\pi_1,\pi_2\}$，其中 $\pi_1=\mathrm{Ind}_B^G(\chi\otimes\chi^{-1})$，$\pi_2=\mathrm{Ind}_B^G(\omega\chi^{-1}\otimes\omega^{-1}\chi)$，$\chi:\Q_p^\times\to k^\times$ 是适当一般的特征。在此背景下有非分裂扩张 $0\to\pi_2\to\kappa\to\pi_1\to0$ 与 $0\to\pi_1\to\kappa'\to\pi_2\to0$（各自在同构意义下唯一），我们将描述各种层 $\mathfrak A(\pi_1)$、$\mathfrak A(\pi_2)$、$\mathfrak A(\kappa)$ 与 $\mathfrak A(\kappa')$。

首先注意，可以写成 $\pi_1=(c\text{-Ind}_K^G\sigma)/(T-\lambda)$ 与 $\pi_2=(c\text{-Ind}_K^G\sigma^{co})/(T-\lambda^{-1})$，其中 $\sigma$ 是某个非常一般的 Serre 权、$\sigma^{co}$ 是它的相伴权，$\lambda$ 适当选取。于是所考虑的所有层，其支撑都包含在 $\pi_\sigma^{-1}(\lambda)$ 中。

轨迹 $\pi_\sigma^{-1}(\lambda)$ 是 $\mathcal U(\sigma|\sigma^{co})_{\mathrm{red}}$ 的闭子叠。沿用 (7.2.23) 的记号，有
\[
\mathcal U(\sigma|\sigma^{co})_{\mathrm{red}}:=[\Spec\bigl(k[t^{\pm1},x,y]/(xy)\bigr)/\mathbb G_m],
\]
其中 $\mathbb G_m$-作用为 $u\cdot(x,y)=(u^2x,u^{-2}y)$。此时映射 $\pi_\sigma:\mathcal U(\sigma|\sigma^{co})_{\mathrm{red}}\to\mathbf A^1$ 就是到 $t$-坐标的投影，故
\[
\pi_\sigma^{-1}(\lambda)=\bigl[(\Spec k[x,y]/(xy))/\mathbb G_m\bigr]
\]
（即 $t=\lambda$ 的轨迹）。

可以给 $k[x,y]/(xy)$ 赋予由 $\mathbb G_m$-作用诱导的分次（于是 $x$ 的次数为 $2$，$y$ 的次数为 $-2$），进而把 $\pi_\sigma^{-1}(\lambda)$ 上的凝聚层描述为分次 $k[x,y]/(xy)$-模。于是对所考虑的层有如下描述：

\ysrc{296}
\[
\begin{aligned}
\mathfrak A(\pi_1)&=\bigl(k[x,y]/(y)\bigr)(-1);\\
\mathfrak A(\pi_2)&=\bigl(k[x,y]/(x)\bigr)(1);\\
\mathfrak A(\kappa)&=\bigl(k[x,y]/(xy)\bigr)(-1);\\
\mathfrak A(\kappa')&=\bigl(k[x,y]/(xy)\bigr)(1).
\end{aligned}
\]
（这里扭照常表示分次的平移。）注意显然的短正合列
\[
0\to\bigl(k[x,y]/(x)\bigr)(1)\xrightarrow{y\cdot}\bigl(k[x,y]/(xy)\bigr)(-1)\to\bigl(k[x,y]/(y)\bigr)(-1)\to0
\]
与
\[
0\to\bigl(k[x,y]/(y)\bigr)(-1)\xrightarrow{x\cdot}\bigl(k[x,y]/(xy)\bigr)(1)\to\bigl(k[x,y]/(x)\bigr)(1)\to0;
\]
它们可以解释为对定义 $\kappa$ 与 $\kappa'$ 的短正合列应用 $\mathfrak A$ 所得。

\subsubsection*{7.4.5.\quad 驯顺主级数型}
在本例与下例中，我们处理某些 $p$-进完备化的紧诱导表示，它们与驯顺潜在 Barsotti--Tate 表示有关。我们满足于只勾勒验证全部断言所需的各项计算。

从驯顺主级数情形开始。为此，设 $\{\sigma,\sigma^{co}\}$ 非常一般，$\tau$ 是驯顺主级数型，使得 $\sigma$ 是 $\overline{\sigma}^{\mathrm{crys},\circ}(\tau)$ 的 Jordan--Hölder 因子。注意这蕴涵 $\sigma^{co}$ 不是 $\overline{\sigma}^{\mathrm{crys},\circ}(\tau)$ 的 Jordan--Hölder 因子；确实，若 $\sigma=\sigma_{a,b}$，则 $\overline{\sigma}^{\mathrm{crys},\circ}(\tau)$ 的 Jordan--Hölder 因子是 $\sigma$ 本身，连同 $\sigma_{a+b,p-1-b}$。（确切地说，型 $\overline{\sigma}^{\mathrm{crys},\circ}(\tau)$ 是从 $\mathrm{GL}_2(\mathbb F_p)$ 的主级数表示膨涨而来的，见 §7.7.7。）

于是可以考虑型 $\tau$ 的潜在 Barsotti--Tate 表示之叠 $\XX^{\mathrm{crys},(1,0),\tau}$。可以相当显式地描述 $\XX$ 的这个闭子叠，下面就来这样做。首先要指出，mod $\varpi$ 纤维 $\XX^{\mathrm{crys},(1,0),\tau}$ 等于 $\XX(\sigma_{a,b})\cup\XX(\sigma_{a+b,p-1-b})$。（这是 §6.1.38 所讨论的几何 Breuil--Mézard 猜想对特定型 $\tau$ 的精确形式。）该并集容许注 7.2.18 给出的显式描述，后者把它嵌入为 $\Spec k[t_1,t_2]/(t_1t_2)\times[\mathbf P^1/\mathbb G_m]$ 的开子叠。后一叠可等同于 $p$-进形式代数叠 $(\Spf(\OO[t_1,t_2]/(t_1t_2-p))^\wedge\times[\widehat{\mathbf P}^1/\widehat{\mathbb G}_m]$ 的底约化子叠；而 $\XX^{\mathrm{crys},(1,0),\tau}$ 同构于该形式代数叠中以 $\XX(\sigma_{a,b})\cup\XX(\sigma_{a+b,p-1-b})$ 为底约化子叠的开子叠。

现在转向描述层 $\mathfrak A\bigl((c\text{-Ind}_K^G\sigma^{\mathrm{crys},\circ}(\tau))^\wedge\bigr)$。当然，这依赖于格 $\sigma^{\mathrm{crys},\circ}(\tau)$ 的具体选取。在所有这些格中，有两个特殊的格 $\Lambda_1$ 与 $\Lambda_2$，由「$\Lambda_1$（相应地 $\Lambda_2$）的余基座（cosocle）等于 $\sigma$（相应地 $\sigma_{a+b,p-1-b}$）」这一条件唯一确定（至缩放）。我们可以且确实把这些格正规化，使 $p\Lambda_1\subset\Lambda_2\subset\Lambda_1$。于是（至缩放）任一格 $\Lambda$ 都形如 $\varpi^a\Lambda_1+\Lambda_2$，其中 $0\le a\le e$（$e$ 为系数域 $L$ 的绝对分歧指数）。（$\mathrm{GL}_2$ 的驯顺主级数型中格的这一描述归功于 Breuil [Bre14, §2]，是 [EGS15, Prop. 4.1.4] 的特例。）

若 $\chi:\mathbb F_p^\times\times\mathbb F_p^\times\to\mathbb Z_p^\times\subseteq\OO^\times$ 记 $\sigma$ 的最高权的驯顺提升，我们把它膨涨为（通常上三角）Iwahori 子群 $I$ 的特征，则 $\Lambda_1=\Ind_I^K\chi$；而若 $\chi^s$ 记 $\chi$ 经非平凡 Weyl 群
\ysrc{297}元素 $s$ 的共轭，则 $\Lambda_2=\Ind_I^K\chi^s$。用 $I$ 在 $G$ 中的正规化子里的元素
\[
\begin{pmatrix}0&1\\p&0\end{pmatrix}
\]
作共轭，交换 $\chi$ 与 $\chi^s$，故有同构
\[
c\text{-Ind}_K^G\Lambda_1=c\text{-Ind}_K^G\chi\cong c\text{-Ind}_I^G\chi^s=c\text{-Ind}_K^G\Lambda_2.
\]
把这个同构与由含入 $\Lambda_2\subset\Lambda_1$ 诱导的含入 $c\text{-Ind}_K^G\Lambda_2\into c\text{-Ind}_K^G\Lambda_1$ 复合，得到 $c\text{-Ind}_K^G\Lambda_1$ 的自同态 $U_p$。交换 $\Lambda_1$ 与 $\Lambda_2$ 的角色，得到 $c\text{-Ind}_K^G\Lambda_1$ 的另一个自同态 $U_p'$，满足 $U_pU_p'=p$；并且
\[
\End_G(c\text{-Ind}_K^G\Lambda_1)=\OO[U_p,U_p']/(U_pU_p'-p).
\]
过渡到 $p$-进完备化，得到
\[
\End_G\bigl((c\text{-Ind}_K^G\Lambda_1)^\wedge\bigr)=\OO[U_p,U_p']^\wedge/(U_pU_p'-p).
\]
若 $\Lambda=\varpi^a\Lambda_1+\Lambda_2$，则 $\Lambda$ 在 $\Lambda_1$ 中的含入诱导等同 $c\text{-Ind}_K^G\Lambda=(\varpi^a,U_p)c\text{-Ind}_K^G\Lambda_1$，从而诱导 $p$-进完备化的相应等同 $(c\text{-Ind}_K^G\Lambda)^\wedge=(\varpi^a,U_p)(c\text{-Ind}_K^G\Lambda_1)^\wedge$。这把任一格 $\Lambda$ 的 $\mathfrak A\bigl((c\text{-Ind}_K^G\Lambda)^\wedge\bigr)$ 之计算化归到 $\mathfrak A\bigl((c\text{-Ind}_K^G\Lambda_1)^\wedge\bigr)$ 的计算。

现在取系数域 $L$ 在 $\Q_p$ 上非分歧（确实，特征 $\chi$、从而格 $\Lambda_1$ 都定义在 $\mathbb Z_p$ 上）；在此假设下看到 $\XX^{\mathrm{crys},(1,0),\tau}$ 正则。于是，由于（如 §7.3.22 所述）层 $\mathfrak A\bigl((c\text{-Ind}_K^G\Lambda_i)^\wedge\bigr)$ 是 Cohen--Macaulay 的，Auslander--Buchsbaum 公式表明它是局部自由模，且事实上必须可逆（例如理由如注 6.1.35 所勾勒）。由于 $\Lambda_1^\vee=\Lambda_2$，且 $(c\text{-Ind}_K^G\Lambda_1)^\wedge\cong(c\text{-Ind}_K^G\Lambda_2)^\wedge$，进一步看出该线丛在 $\XX^{\mathrm{crys},(1,0),\tau}$ 上是 Grothendieck--Serre 自对偶的。

现在可以验证 $\mathrm{Pic}([\widehat{\mathbf P}^1/\widehat{\mathbb G}_m])=\mathbf Z\times\mathbf Z$，其中第一个 $\mathbf Z$ 来自 $\mathrm{Pic}(\widehat{\mathbf P}^1)=\mathbf Z$，第二个 $\mathbf Z$ 来自按 $\widehat{\mathbb G}_m$ 的特征作扭。以 $\OO(m)_{\widehat{\mathbf P}^1}(n)$ 记对应于配对 $(m,n)\in\mathbf Z\times\mathbf Z$ 的可逆层。沿复合
\[
\begin{aligned}
f:\XX^{\mathrm{crys},(1,0),\tau}\;&\xrightarrow{\ \text{open}\ }(\Spf(\OO[t_1,t_2]/(t_1t_2-p))^\wedge\times[\widehat{\mathbf P}^1/\widehat{\mathbb G}_m]\\
&\xrightarrow{\ \text{proj.}\ }[\widehat{\mathbf P}^1/\widehat{\mathbb G}_m]
\end{aligned}
\]
拉回，诱导同构
\[
\mathrm{Pic}([\widehat{\mathbf P}^1/\widehat{\mathbb G}_m])\isom\mathrm{Pic}(\XX^{\mathrm{crys},(1,0),\tau}).
\]
此外，$[\widehat{\mathbf P}^1/\widehat{\mathbb G}_m]$ 的对偶层，即 $\OO(-2)_{\widehat{\mathbf P}^1}$，拉回到 $\XX^{\mathrm{crys},(1,0),\tau}$ 的对偶层。于是（由于它自对偶）发现
\[
\mathfrak A\bigl((c\text{-Ind}_K^G\Lambda_1)^\wedge\bigr)\cong f^*\OO(-1)_{\widehat{\mathbf P}^1}.
\]
由于 $\mathfrak A$ 全忠实，还得到同构
\[
\begin{aligned}
\End_G\bigl((c\text{-Ind}_K^G\Lambda_1)^\wedge\bigr)&\isom\Gamma\bigl(\XX^{\mathrm{crys},(1,0),\tau},\OO_{\XX^{\mathrm{crys},(1,0),\tau}}\bigr)\\
&=\OO[t_1,t_2]^\wedge/(t_1t_2-p),
\end{aligned}
\]
它（适当选取 $t_i$ 后）由 $U_p\mapsto t_1$、$U_p'\mapsto t_2$ 给出。对任一格 $\Lambda=\varpi^a\Lambda_1+\Lambda_2$，随后发现
\[
\mathfrak A\bigl((c\text{-Ind}_K^G\Lambda)^\wedge\bigr)=(\varpi^a,t_1)f^*\OO_{\widehat{\mathbf P}^1}(-1).
\]

\ysrc{298}
为了把这些 $\mathfrak A$ 的取值与上面为 Serre 权的紧诱导所算的取值作比较，先注意拉回到特殊纤维诱导同构
\[
\mathrm{Pic}([\widehat{\mathbf P}^1/\widehat{\mathbb G}_m])\isom\mathrm{Pic}([\mathbf P^1/\mathbb G_m]).
\]
然后，若以 $x_0,x_1$ 记 $\mathbf P^1$ 的齐次坐标（于是注 7.2.18 的非齐次坐标 $x$ 等于 $x_1/x_0$），以 $u$ 记 $\mathbb G_m$ 上的坐标，则 $[\mathbf A^1_{x_1/x_0}/(\mathbb G_m)_u]$ 与 $[\mathbf A^1_{x_0/x_1}/(\mathbb G_m)_{u^{-1}}]$ 各自嵌入为 $[\mathbf P^1/\mathbb G_m]$ 的开子叠；限制到这些子叠诱导嵌入
\[
\begin{aligned}
\mathrm{Pic}([\mathbf P^1/\mathbb G_m])&\into\mathrm{Pic}([\mathbf A^1_{x_1/x_0}/(\mathbb G_m)_u])\times\mathrm{Pic}([\mathbf A^1_{x_0/x_1}/(\mathbb G_m)_{u^{-1}}])\\
&=\mathbf Z\times\mathbf Z,
\end{aligned}
\]
由 $(m,n)\mapsto(m+n,m-n)$ 给出。由此计算限制
\[
(f^*\OO(-1)_{\widehat{\mathbf P}^1})|_{\XX(\sigma_{a,b})}=\OO_{\XX(\sigma_{a,b})}(-1)
\]
与
\[
(f^*\OO(-1)_{\widehat{\mathbf P}^1})|_{\XX(\sigma_{a+b,p-1-b})}=\OO_{\XX(\sigma_{a+b,p-1-b})}(-1).
\]
于是，若把 $\mathfrak A$ 作用于短正合列
\[
0\to(c\text{-Ind}_K^G\Lambda_1)^\wedge\xrightarrow{U_p}(c\text{-Ind}_K^G\Lambda_1)^\wedge\to c\text{-Ind}_K^G\sigma_{a,b}\to0,
\]
相应地
\[
0\to(c\text{-Ind}_K^G\Lambda_1)^\wedge\xrightarrow{U_p'}(c\text{-Ind}_K^G\Lambda_1)^\wedge\to c\text{-Ind}_K^G\sigma_{a+b,p-1-b}\to0,
\]
便得到层的短正合列
\[
0\to f^*\OO_{\widehat{\mathbf P}^1}(-1)\xrightarrow{t_1}f^*\OO_{\widehat{\mathbf P}^1}(-1)\to\OO_{\XX(\sigma_{a,b})}(-1)\to0,
\]
相应地
\[
0\to f^*\OO_{\widehat{\mathbf P}^1}(-1)\xrightarrow{t_2}f^*\OO_{\widehat{\mathbf P}^1}(-1)\to\OO_{\XX(\sigma_{a+b,p-1-b})}(-1)\to0.
\]

\subsubsection*{7.4.6.\quad 驯顺尖点型}
对驯顺尖点型的情形，可以作与 §7.4.5 类似的分析。即，存在驯顺尖点型 $\tau$，由「$\sigma^{\mathrm{crys},\circ}(\tau)$ 的 Jordan--Hölder 因子是 $\sigma$ 与 $\sigma^{co}$」这一条件唯一确定。此时潜在 Barsotti--Tate 子叠 $\XX^{\mathrm{crys},(1,0),\tau}$ 是 $\XX(\sigma|\sigma^{co}):=\XX(\sigma)\cup\XX(\sigma^{co})$ 的 $p$-进形式提升。如注 7.2.17 所述，$\XX(\sigma|\sigma^{co})$ 嵌入为 $\mathbf P^1\times[(\Spec k[x,y]/(xy))/\mathbb G_m]$ 的开子叠；而形变论论证表明 $\XX^{\mathrm{crys},(1,0),\tau}$ 同构于 $p$-进形式代数叠 $\widehat{\mathbf P}^1\times[(\Spf(\OO[x,y]/(xy-p))^\wedge/\widehat{\mathbb G}_m]$ 的提升 $\XX(\sigma|\sigma^{co})$ 的开形式代数子叠。\tn{底本此处同样作「$[(\Spf(\OO[x,y]/(xy-p))^\wedge/\widehat{\mathbb G}_m]$」，左圆括号无配对；译文照录。}

同主级数情形一样，$\sigma^{\mathrm{crys}}(\tau)$ 中有两个特殊的格 $\Lambda_1$ 与 $\Lambda_2$，各自至缩放唯一确定，其余基座分别为 $\sigma$ 与 $\sigma^{co}$。固定 $\Lambda_1$ 的选取后，可以缩放 $\Lambda_2$ 使 $p\Lambda_1\subset\Lambda_2\subset\Lambda_1$；任一格都（至缩放）形如 $\Lambda=\varpi^a\Lambda_1+\Lambda_2$，其中 $0\le a\le e$。于是为描述各种层 $\mathfrak A\bigl((c\text{-Ind}_K^G\Lambda)^\wedge\bigr)$，只需描述层 $\mathfrak A\bigl((c\text{-Ind}_K^G\Lambda_1)^\wedge\bigr)$ 与 $\mathfrak A\bigl((c\text{-Ind}_K^G\Lambda_2)^\wedge\bigr)$，以及由含入 $\Lambda_2\subset\Lambda_1$ 诱导的态射
\begin{equation}\tag{7.4.7}\label{eq:a4-7-4-7}
\mathfrak A\bigl((c\text{-Ind}_K^G\Lambda_2)^\wedge\bigr)\to\mathfrak A\bigl((c\text{-Ind}_K^G\Lambda_1)^\wedge\bigr).
\end{equation}

% 本片止于印刷页 298 页末（§7.4.6，公式 (7.4.7)），续片自 PDF 312 页（印刷 299）起
